\documentclass[instructions.tex]{subfiles}
\begin{document}

\chapter{On ne doit pas différer d’un jour sa conversion.}

Le temps présent et l’avenir doivent vous engager à vous convertir incessamment.

\primo Le temps présent. Vous êtes dans le péché et dans le danger d’y mourir. Dieu vous offre aujourd’hui la grâce pour en sortir ; pourquoi ne le faites-vous pas ? Dire que rien ne presse, c’est raisonner en insensé. Attend-on au lendemain pour ôter les taches que l’on aperçoit sur son visage ? Quand on est dangereusement blessé, diffère-t-on d’appliquer le remède ? Quand on est attaqué par une bête féroce, ne crie-t-on pas d’abord au secours ? Pourquoi n’agissez-vous pas de même, lorsque votre âme est souillée par le péché, qu’elle est blessée, qu’elle est au pouvoir du démon ? Pourquoi ne pas implorer aussitôt le secours du ciel, et recourir à la pénitence ?

Vous ne voudriez pas mourir dans le péché ; vous espérez donc qu’un jour vous le quitterez ; mais pourquoi ne le quittez-vous pas aujourd’hui, puisque vous en avez le temps ? Croyez-vous que dans la suite vous serez mieux disposé à le faire ? Dieu sera-t-il demain plus disposé à vous pardonner ? Les mauvaises habitudes seront-elles plus faciles à surmonter ? Votre cœur sera-t-il moins endurci ? Loin de là, votre conversion deviendra toujours plus difficile ; les mauvaises habitudes, toujours plus fortes. Le temps, qui affaiblit tout, ne peut les empêcher de croître.

\secundo Quant à l’avenir, les desseins de Dieu sont si redoutables, que vous ne sauriez différer une heure votre conversion, sans vous mettre en danger d’être perdu. Savez-vous le nombre des péchés que Dieu veut souffrir de vous, et la mesure des grâces qu’il veut vous accorder ? Savez-vous jusqu’où doit aller sa patience ? Ne craignez-vous point que le premier péché mortel que vous commettrez, ne soit le dernier que Dieu veuille souffrir de vous ?

Tous les malfaiteurs ne sont pas également punis ; les uns vieillissent dans le brigandage, les autres sont surpris dès le premier crime, et punis de mort. Combien d’anges et de millions d’hommes n’ont commis qu’un seul péché ! Ce péché seul a suffi pour consommer leur réprobation. Vous avez donc échappé à l’enfer, autant de moments que vous avez vécu dans le péché mortel ; et vous êtes toujours exposé à y tomber, pendant que vous vivez dans cet état.

Un aveugle qui marche vers un précipice, y tombe enfin. Le dernier pas qu’il fait n’est pas plus grand que les autres ; c’est cependant ce dernier pas qui le fait tomber. Pour tomber dans l’enfer, il n’est pas nécessaire de commettre de plus grands crimes ; c’est assez de continuer dans l’état où vous êtes. Voici long-temps que vous y persévérez, que vous avancez vers l’abîme ; vous n’en êtes éloigné que d’un pas ; peut-être cette nuit on demandera votre âme. À quoi vous exposez-vous en différant votre conversion ?

Vous espérez que vous aurez du temps à l’avenir. Vous êtes un téméraire qui espérez ce que Dieu ne vous a point promis. Il vous promet le pardon si vous retournez sincèrement à lui ; mais il ne vous a pas promis un jour, pas même une heure. Il vous menace au contraire de vous enlever de ce monde à l’heure que vous y penserez le moins\footnote{Quâ horâ non putatis, Luc. 12. 40.}. Peut-être serez-vous demain en vie, mais peut-être aussi n’y serez-vous plus. N’êtes-vous pas aveugle de fonder votre salut sur un peut-être ? Mais, quand vous seriez assuré d’être demain en vie, êtes-vous certain que vous aurez la grâce ? Dieu est-il obligé d’attendre que vous soyez d’humeur à le recevoir ? et méritez-vous qu’il vous l’accorde, après avoir abusé de tant de grâces, et négligé par votre malice tant d’occasions de salut ? Dieu vous donnera-t-il ces grâces victorieuses dont il favorise ses serviteurs ? Dieu le peut, il les a même accordées à de grands pécheurs ; mais devez-vous vous y attendre ?

S’il n’est point de faveur plus grande que la grâce efficace, aussi n’est-il point d’espérance plus téméraire que celle du pécheur qui s’imagine qu’il ne manquera pas de l’avoir quand il voudra. Doit-on s’attendre que Dieu donne des grâces extraordinaires à celui qui abuse de toutes les grâces ordinaires ? Vous en avez eu de ces grâces puissantes, vous en avez abusé. Lorsque Dieu vous en donnera, n’en abuserez-vous pas encore ? celui qui néglige de se corriger, devient de jour en jour plus endurci et plus incorrigible.

Le délai de la conversion est de tous les pièges le plus dangereux. Le démon n’a garde de persuader à un pécheur qu’il ne faut pas se convertir ; mais le démon est content s’il peut l’engager à différer. Quand on retarde de quelques jours sa conversion, on la remettra bientôt à un mois, ensuite à une année, enfin à la mort. Oh ! qu’il est affligeant de s’apercevoir des pièges du démon quand on ne peut plus les éviter, et de sentir les forces de son ennemi quand on est affaibli et qu’on se meurt !

\section{Résolutions}

Pourrais-je penser sans frémir à quels dangers un plus long délai de conversion m’exposerait ? 

\primo Aujourd’hui donc, à l’heure même, je commence. 

\secundo Oui, ce jour ne s’écoulera pas sans que je me sois concerté avec mon sage guide spirituel sur les mesures que je dois prendre pour restituer ce bien mal acquis, pour rompre cette liaison dangereuse, pour vaincre cette tentation séduisante, pour dompter ce penchant véhément qui est en moi la cause de tant de péchés. 

\prayer{Mon Dieu, aidez-moi à sortir de l’état de fausse sécurité dans lequel j’ai vécu jusqu’à présent.}

\end{document}
